%%%%%%%%%%%%%%%%%%%%%%%%%%%%%%%%%%%%%%%%%%%%%%%%%%%%%%%%%%%%%%%%%%%%%%%%
\chapter*{Introduzione}
%%%%%%%%%%%%%%%%%%%%%%%%%%%%%%%%%%%%%%%%%%%%%%%%%%%%%%%%%%%%%%%%%%%%%%%%

I tensori sono oggetti che modellano molti esperimenti in un numero disparato di aree scientifiche.
\todo{scrivi altro}
Decomporli in oggetti più facili da analizzare \ldots \todo{scrivi altro}
Studieremo la \emph{canonical polyadic decomposition}, nota anche come PARAFAC, per ``parallel factor
analysis'', CANDECOMP, per ``canonical decomposition'', e
CP \emph{decomposition}, il quale nome potrebbe o come acronimo della combinazione di CANDECOMP e
PARAFAC.

Le decomposizioni di tensori occorrono in numerose aree scientifiche. Alcune di queste sono:
nell'ambito medico della psicometria è stata impiegata per studiare gli impulsi celebrali [311,
		71, 158], in particolare nel trovare l'area del cervello di un paziente in cui è in atto una crisi
epilettica, o la diagnosi e gestione di ictus.; sempre nella medicina,
nell'interpretare la immagini di risonanze magnetiche (MRI); vedi ad esempio e relative
reference.
Nella geofisica, ad esempio l'interpretazione di data magnetotellurica per strutture
regionali a una e due dimensioni, ad esempio in \ldots e le reference all'interno (wtf?).
Nell'analisi numerica.\todo{questo va o detto a parole, oppure descrivere esempio delle funzioni
	semiseparabili} Grazie alla convergenza delle serie di Fourier nello spazio delle
funzioni $ f : \R^n \to \R $ misurabili secondo la norma euclidea $ L^2(\R^n) $, si
ha $ L^2(\R^n) =
	L^2(\R) \otimes \cdots \otimes L^2(\R) $, e spesso vengono approssimate funzioni di $ n $
variabili da una somma finita di prodotti di funzioni in una variabile, generalizzando la
separazione di variabili. Si veda \ldots

Questo elaborato si sviluppa con la seguente struttura:

Nel capitolo 1, introdurremo i tensori e i relativi spazi. Definiremo la fattorizzazione CP e del
relativi metodi di calcolo. Discuteremo alcune loro caratteristiche
fondamentali che li differenziano dalle matrici.

Nel capitolo 2 dimostreremo il Teorema di Kruskal, che afferma l'unicità della scrittura della
fattorizzazione CP di un tensore sotto opportune ipotesi, discutendo alcuni casi di non unicità.

Nel capitolo 3 sposteremo il focus allo studio della complessità della moltiplicazione matriciale,
studiando gli \emph{algortmi di moltiplicazione veloci} dal punto di vista della decomposizione
CP.

Nel capitolo 4 presenteremo alcune applicazioni modellistiche della CP, presentando due esperimenti
numerici: il recupero della concentrazione di sostanze chimiche in un composto (\emph{spettroscopia
	fluorescente}) e l'analisi di segnali (\emph{cocktail party problem}).

Le principali fonti di questo lavoro sono TENSORB e LANDSBERG1,2.

Come descritto da Landsberg\todo{cite}, nell'ambito del calcolo tensoriale avviene uno ``scontro di
culture'', dove confluiscono risultati teorici trovati dagli studiosi delle applicazioni scientifiche
(e.g., psicometria, chemiometria, neurologia), analisti numerici, teorici della complessità e
geometri algebrici.

Credits: cite ballard E Kolda per aver inspirato molte figure. Librerie numeriche.
